% ==============================================================================
% config/metadados.tex
% Metadados institucionais centralizados
% ==============================================================================

\titulo{Arquitetura e Especificação Técnica de uma Pipeline Híbrida de Análise de Dados Industriais Integrando WEG PLC500, Comunicação Modbus TCP/IP e Orquestração com Node-RED}

\autor{Augusto P. Dantas \and Gabriel Simão \and Gabriel Silvano}
\newcommand{\autorespdf}{Augusto P. Dantas, Gabriel Simão, Gabriel Silvano}

\local{Jaraguá do Sul -- SC}

\data{2026}

\instituicao{Centroweg Senai}

\tipotrabalho{Relatório Técnico}

\orientador[Coordenador:]{Lucas de Sousa Santos}

% --- Identificação técnica do documento ---
\newcommand{\codigorelatorio}{RT-AUT-2026-001}
\newcommand{\versaorelatorio}{Revisão 1.0}
\newcommand{\classificacaodoc}{Documentação de Engenharia e Arquitetura}

\preambulo{%
  Relatório técnico desenvolvido no escopo de Situação de Aprendizagem de
  Internet das Coisas (IoT) no Centroweg Senai, apresentando a arquitetura, a
  especificação, a implementação completa e a validação por simulação em gêmeo
  digital de uma pipeline híbrida de análise de dados industriais integrando
  controlador lógico programável WEG PLC500 programado em \textit{Structured Text}
  (IEC~61131-3) via CODESYS, comunicação Modbus TCP/IP, orquestração em cinco
  abas no Node-RED, inferência de aprendizado de máquina via FastAPI e
  \textit{scikit-learn}, diagnóstico por modelo de linguagem restrito por
  JSON Schema Draft 2020-12 e persistência relacional em PostgreSQL.%
}
